\noindent\textbf{Por qué se usa este procedimiento.} La tracción es $\boldsymbol t=\boldsymbol\sigma\boldsymbol n$. En tracción uniaxial, proyectar sobre la normal da $\sigma_n=\sigma\cos^2\theta$. Hay dos proyecciones: una forma el vector tracción y otra su componente normal.

\begin{enumerate}

\item $\cos 30^{\circ} = \sqrt3/2$, de modo que $\cos^2 30^{\circ} = 3/4$.

\end{enumerate}

\noindent\textbf{Operaciones y resultados intermedios.}

\begin{enumerate}

\item Tensión normal.

\[\fitmath{y_{1}=100\cdot {\cos\left(\frac{\pi}{6}\right)}^{2}\approx 75\,\mathrm{MPa}}\]

\end{enumerate}
